\section*{第 \ref{ch:g11:infrared} 章——红外光谱}

\begin{solution}{exo:g11:infrared:1}
$\lambda = 1/3300 = \qty{3.03e-4}{cm} = \qty{3.03}{\micro m}$;
$\lambda = 1/1050 = \qty{9.52e-4}{cm} = \qty{9.52}{\micro m}$。
\qty{3300}{cm^{-1}} 处的吸收带，\omterm{def:g11:infrared:wavenumber}{波数}较大，属于较快的
振动。
\end{solution}

\begin{solution}{exo:g11:infrared:2}
$\qty{10}{\micro m} = \qty{1.0e-3}{cm}$，故 $\sigma = 1/\num{1.0e-3} =
\qty{1000}{cm^{-1}}$.
\end{solution}

\begin{solution}{exo:g11:infrared:3}
酮的 \ce{C=O}；\omterm{def:g11:polarity-and-cohesion:hydrogen-bond}{氢键}缔合醇的 \ce{O-H}；
碳链上的 \ce{C-H} 键。
\end{solution}

\begin{solution}{exo:g11:infrared:4}
\qty{1715}{cm^{-1}} 处的吸收带：\ce{C=O} 键。
\end{solution}

\begin{solution}{exo:g11:infrared:5}
仪器测量穿过样品的光；
某键吸收光时，穿过的光减少，曲线下凹。$T =
\qty{100}{\percent}$ 表示该\omterm{def:g11:infrared:wavenumber}{波数}的光均未被
吸收。
\end{solution}

\begin{solution}{exo:g11:infrared:6}
第一种物质在 \qty{1715}{cm^{-1}} 处有 \ce{C=O} 吸收带，无 \ce{O-H} 吸收带。
它是酮：\ce{CH3-CO-CH2-CH3}，丁酮（又称
丁-2-酮）。第二种有 \ce{O-H} 吸收带，无 \ce{C=O}：
为醇，例如丁-1-醇。
\end{solution}

\begin{solution}{exo:g11:infrared:7}
酸的 \ce{O-H} 基团\omterm{def:g11:polarity-and-cohesion:hydrogen-bond}{氢键}作用很强（
\omterm{def:g7:atoms-and-molecules:molecule}{分子}两两配对，每个 \ce{O-H} 与另一\omterm{def:g7:atoms-and-molecules:molecule}{分子}的 \ce{C=O} 成键），
且每个\omterm{def:g7:atoms-and-molecules:molecule}{分子}的成键情况各不相同，吸收带
进一步弥散，分布在很宽的范围内。
\end{solution}

\begin{solution}{exo:g11:infrared:8}
在液体（A）中，\ce{O-H} 吸收带较宽，
位于 \qty{3350}{cm^{-1}} 附近；在气体（B）中，\omterm{def:g7:atoms-and-molecules:molecule}{分子}相距远，
不形成\omterm{def:g11:polarity-and-cohesion:hydrogen-bond}{氢键}，吸收带尖锐，
位于 \qty{3666}{cm^{-1}} 附近。
\end{solution}

\begin{solution}{exo:g11:infrared:9}
很难判断：醛在 \qty{1730}{cm^{-1}} 附近，
酮在 \qty{1715}{cm^{-1}} 附近，两者数值接近。将整个
光谱（含\omterm{def:g11:infrared:band}{指纹区}）与参考光谱比较，
或借助其他技术才能确定。
\end{solution}

\begin{solution}{exo:g11:infrared:10}
酯（\ce{C=O} 在 \qty{1735}{cm^{-1}} 附近，\ce{C-O} 吸收强，
无 \ce{O-H}），例如乙酸乙酯。乙酸的\omterm{def:g11:organic-skeletons:formulas}{分子式}为 \ce{C2H4O2}，
而不是 \ce{C4H8O2}。丁酸的\omterm{def:g11:organic-skeletons:formulas}{分子式}虽符合，但会显示
酸特有的很宽 \ce{O-H} 吸收带，应予排除。
\end{solution}

\begin{solution}{exo:g11:infrared:11}
丙-1-醇是醇，与 A 相同。丙酸是\omterm{def:g11:functional-groups:carboxylic-acid}{羧酸}：
与 D 相同。
\end{solution}

\begin{solution}{exo:g11:infrared:12}
\qty{3350}{cm^{-1}} 附近较宽的 \ce{O-H} 吸收带减弱，
而位于 \qty{1715}{cm^{-1}} 附近的强 \ce{C=O} 吸收带增强。
反应完成的标志是 \ce{O-H} 吸收带消失，
\ce{C=O} 吸收带不再增强。
\end{solution}

\begin{solution}{exo:g11:infrared:13}
丁酸的 \ce{O-H} 吸收带很宽，范围
为 2500 至 \qty{3300}{cm^{-1}}，酯则无；
酸的 \ce{C=O} 吸收峰位于 \qty{1710}{cm^{-1}} 附近，较宽；酯的位于 \qty{1735}{cm^{-1}} 附近，
较尖锐，在 \qty{1240}{cm^{-1}} 附近有强 \ce{C-O} 吸收。
酸中的 \ce{O-H} 基团形成\omterm{def:g11:polarity-and-cohesion:hydrogen-bond}{氢键}，
这也使其 \ce{C=O} 略有减弱。
\end{solution}

\begin{solution}{exo:g11:infrared:14}
残留的乙醇（\omterm{def:g10:synthesis-yield:synthesis}{合成}该酯所用的醇）或水（来自
\omterm{def:g4:air-a-mixture-of-gases:air}{空气}或洗涤过程）：两者均含\omterm{def:g11:polarity-and-cohesion:hydrogen-bond}{氢键}缔合的 \ce{O-H} 基团。
样品可与参考光谱比较，或经干燥、\omterm{def:g10:chemical-species:distillation}{蒸馏}，
或测量其沸点。
\end{solution}

\begin{solution}{exo:g11:infrared:15}
同样两\omterm{def:g7:atoms-and-molecules:atom}{原子}间的\omterm{def:g10:lewis-and-shape:multiple-bond}{双键}相当于劲度系数更大的弹簧，
因此振动更快，\omterm{def:g11:infrared:wavenumber}{波数}更高（1715 对 \qty{1050}{cm^{-1}}）。
\ce{O-H}、\ce{N-H} 和 \ce{C-H} 键都含有氢\omterm{def:g7:atoms-and-molecules:atom}{原子}，
而氢是最轻的\omterm{def:g7:atoms-and-molecules:atom}{原子}。轻球挂在弹簧上，振动得快，
因此\omterm{def:g11:infrared:wavenumber}{波数}在 2900--\qty{3600}{cm^{-1}} 附近。
\end{solution}

\begin{solution}{pb:g11:infrared:1}
\textbf{1.} \ce{CH3-CH2-OH}; \ce{CH3-CO-CH3}; \ce{CH3-COOH}.

\textbf{2.} \ce{C2H6O}; \ce{C3H6O}; \ce{C2H4O2}.

\textbf{3.} 羟基；羰基（酮）；羧基。

\textbf{4.} 乙醇：\ce{O-H} 和 \ce{C-H}。丙酮：
\ce{C-H} 和 \ce{C=O}。乙酸：\ce{O-H}、\ce{C-H} 和 \ce{C=O}。

\textbf{5.} 无 \ce{C=O} 吸收带，有宽的 \ce{O-H} 吸收带：乙醇。

\textbf{6.} 在 2500 至 \qty{3300}{cm^{-1}} 之间有很宽的 \ce{O-H} 吸收带，
且在 \qty{1710}{cm^{-1}} 附近有强的 \ce{C=O} 吸收带：乙酸。

\textbf{7.} 丙酮：在
\qty{1715}{cm^{-1}} 处有很强的 \ce{C=O} 吸收带，无 \ce{O-H}。

\textbf{8.} 可能是（醋、\omterm{def:g6:solutions-and-solubility:solute}{溶剂}、醇），但随意闻未知
液体有危险；而光谱既安全又可靠。

\textbf{9.} 酸的 \ce{O-H} 基团形成的\omterm{def:g11:polarity-and-cohesion:hydrogen-bond}{氢键}更强
（\omterm{def:g7:atoms-and-molecules:molecule}{分子}两两缔合），因此吸收带位置更低，且宽得多。

\textbf{10.} \ce{CH3-O-CH3}，\ce{C2H6O}；不含 \ce{O-H}，因此
无 \ce{O-H} 吸收带。

\textbf{11.} \ce{C=O} 吸收带：丙醛在 \qty{1730}{cm^{-1}}
附近，丙酮在 \qty{1715}{cm^{-1}} 附近。

\textbf{12.} \ce{O-H} 吸收带：酯不含 \ce{O-H}。其 \ce{C=O}
吸收带应位于 \qty{1735}{cm^{-1}} 附近。

\textbf{13.} \omterm{def:g11:organic-skeletons:isomer}{同分异构体}的\omterm{def:g11:organic-skeletons:formulas}{分子式}相同；其\omterm{def:g11:functional-groups:functional-group}{官能团}不同，
而光谱能显示出这些\omterm{def:g11:functional-groups:functional-group}{官能团}。

\textbf{14.} \qty{1715}{cm^{-1}}.

\textbf{15.} $1715 \times 100 = \qty{1.715e5}{m^{-1}}$.

\textbf{16.} $\lambda = 1/\num{1.715e5} = \qty{5.83e-6}{m} =
\qty{5.83}{\micro m}$.

\textbf{17.} 否：可见光的范围大约从 0.4 到
\qty{0.75}{\micro m}。它是红外线。

\textbf{18.} $\num{1.05e5}\ \si{m^{-1}}$，故 $\lambda = \qty{9.52e-6}{m}
= \qty{9.52}{\micro m}$.

\textbf{19.} \ce{C=O} 键（\omterm{def:g11:infrared:wavenumber}{波数}较大，波长较短）。

\textbf{20.} 约为 \qty{5.8}{\micro m}。
\end{solution}
